\section{Fundamentação Teórica}
\label{sec:fundamentacao}
Esta seção apresenta os conceitos fundamentais necessários para a compreensão do problema investigado neste trabalho. Inicialmente, são discutidos os Modelos de LLMs e sua evolução para agentes capazes de interagir com sistemas externos. Em seguida, é apresentado o MCP, protocolo responsável por padronizar a comunicação entre agentes e ferramentas. Posteriormente, são abordadas as descrições em linguagem natural utilizadas por ferramentas MCP, elemento central para o processo de tomada de decisão dos agentes. Por fim, são discutidos os conceitos relacionados aos \textit{smells} em descrições de ferramentas, bem como seu impacto na qualidade das interações realizadas por agentes baseados em LLMs.

\subsection{Modelos de Linguagem de Grande Escala}

A \textbf{Inteligência Artificial (IA)} refere-se ao campo que estuda e desenvolve sistemas capazes de exibir inteligência baseada em máquinas em oposição à biológica \cite{Taraghi2026}. Os \textbf{LLMs} constituem um avanço transformador dentro deste campo, situando-se no subdomínio da \textbf{IA generativa} \cite{mcpAtFirstGlance, multiagent}. Enquanto a IA é uma área ampla, os LLMs são modelos de fundação especializados no processamento de linguagem natural por meio de \textbf{probabilidade e raciocínio dinâmico} \cite{2025arXiv250814704L}. Estes modelos são treinados sobre \textbf{extenso acervo estruturado de textos}, o que os torna aptos a compreender, gerar e manipular linguagem em múltiplos contextos \cite{mcpAtFirstGlance, NEURIPS2023_91f18a12}.

As funcionalidades desses modelos abrangem tarefas diversificadas, incluindo tradução, sumarização e raciocínio assistido \cite{2025arXiv250814704L, 11298840}. O amadurecimento dessas capacidades impulsionou o surgimento de \textbf{agentes inteligentes} projetados para decompor problemas complexos em subtarefas gerenciáveis \cite{2025arXiv250814704L}. Tais agentes operam de forma \textbf{parcialmente autônoma}, executando fluxos de trabalho com independência operacional, embora ainda exijam supervisão humana e validação estratégica \cite{multiagent}. Essa transição do modelo como simples gerador de texto para orquestrador de recursos externos amplia significativamente seu potencial de aplicação em cenários reais que exigem dados atualizados ou ações no mundo físico \cite{2025arXiv250814704L}.



Apesar de seu conhecimento ser restrito às informações disponíveis durante o treinamento, os LLMs apresentam limitações inerentes decorrentes dessa natureza estática. Para superar tais restrições, surgiu o paradigma de \textit{tool calling}, por meio do qual agentes baseados em LLMs podem invocar ferramentas externas, acessar bancos de dados, consultar Interfaces de Programação de Aplicação 
(\textit{Application Programming Interface} - API) e interagir com sistemas computacionais diversos \cite{mcpAtFirstGlance}.Dessa forma, os modelos deixam de atuar apenas como geradores de texto e passam a desempenhar o papel de orquestradores de recursos externos, transformando-se de sistemas estáticos em centros de coordenação dinâmica. Nesse papel de orquestrador, o LLM deixa de apenas prever a próxima palavra e passa a interpretar a intenção do usuário para formular requisições estruturadas que podem ser compreendidas por sistemas externos \cite{2025arXiv250814704L}.

\subsection{Model Context Protocol}

O \textit{Model Context Protocol} (MCP) foi proposto com o objetivo de padronizar a integração entre LLMs e fontes externas de informação, reduzindo a complexidade envolvida na comunicação entre agentes de IA e sistemas heterogêneos \cite{shi2026description}. O protocolo adota uma arquitetura cliente-servidor na qual agentes atuam como clientes e se conectam a servidores MCP responsáveis por disponibilizar recursos, dados e ferramentas de forma padronizada.

Essa abordagem busca fornecer um mecanismo uniforme de integração entre IAs e sistemas externos, independentemente da tecnologia empregada pelos sistemas. De maneira análoga ao papel desempenhado pelo protocolo HTTP na \textit{Web}, o MCP estabelece convenções para descoberta, descrição e utilização de recursos, promovendo interoperabilidade entre diferentes aplicações \cite{guo2025measurement}. Como consequência, desenvolvedores podem disponibilizar funcionalidades para agentes inteligentes sem a necessidade de implementar integrações específicas para cada plataforma ou modelo de IA utilizado.

\subsection{Descrições de Ferramentas MCP}
No ecossistema MCP, as ferramentas (\textit{tools}) representam funcionalidades invocadas por agentes para realizar ações ou coletar dados externos. Em sistemas tradicionais, a interação ocorre via \textbf{assinaturas de código} (o nome técnico da função) e \textbf{mecanismos de tipagem} (regras estritas sobre o formato e tipo dos dados), exigindo chamadas precisas. Entretanto, agentes baseados em LLMs dependem de descrições em linguagem natural para identificar as \textbf{capacidades} disponíveis, ou seja, as tarefas específicas que a ferramenta consegue executar \cite{2026arXiv260214878M}. Essas documentações detalham o propósito, as restrições e o significado dos parâmetros, funcionando como a principal interface de comunicação. Assim, a qualidade desses textos determina a precisão do agente ao selecionar e empregar os recursos técnicos necessários para cumprir seus objetivos.

Uma descrição de ferramenta normalmente contém informações relacionadas ao propósito da funcionalidade, aos cenários adequados de utilização, às restrições de uso e ao significado de seus parâmetros. A partir dessas informações, o agente de IA decide quando utilizar determinada ferramenta e quais argumentos devem ser fornecidos durante sua execução \cite{2026arXiv260214878M}. Consequentemente, a qualidade da descrição da ferramenta influencia na capacidade do agente de utilizar corretamente as ferramentas e empregá-las de maneira coerente com seus objetivos.

\subsection{Smells em Descrições de Ferramentas}
O termo \textit{smell} é utilizado na Engenharia de Software para designar características que sugerem potenciais problemas de qualidade em um artefato, mesmo quando não configuram defeitos funcionais. Esses indícios podem dificultar atividades de compreensão, manutenção e evolução de sistemas, motivando o desenvolvimento de diferentes técnicas de detecção ao longo dos anos. O uso recente de LLMs para identificar \textit{smells}, demonstrando que esses modelos são capazes de alcançar resultados iguais ou superiores aos obtidos por abordagens tradicionais baseadas em regras ou métricas estáticas \cite{souza2026beyond}. 

No contexto do MCP, o conceito de \textit{smell} é aplicado às descrições em linguagem natural que acompanham as ferramentas disponibilizadas aos agentes. Os \textit{smells} correspondem a ambiguidades, omissões, inconsistências ou outros problemas estruturais capazes de comprometer a compreensão adequada da ferramenta por parte do modelo de IA \cite{2026arXiv260214878M}. Para avaliar a qualidade dessas descrições, a literatura propõe rubricas baseadas em critérios como clareza do propósito, completude das instruções e detalhamento dos parâmetros. A presença desses problemas pode levar agentes a selecionar ferramentas inadequadas, tornando a detecção de \textit{smells} um aspecto fundamental para a confiabilidade de ecossistemas baseados em MCP \cite{2026arXiv260214878M}.